\documentclass[12pt]{article}
\usepackage[ left=0.5in,right=0.5in ]{geometry}
\usepackage{../styles/ntalgo}

\usepackage[backend=biber,style=alphabetic]{biblatex}
\addbibresource{refs.bib}

\begin{document}

\title{Low Diameter Decompositions}
\author{Navid Tajkhorshid}

\maketitle

I'm writing this mainly as a review for myself. I'm trying to do everything from
memory/rederive things, so the results here may not be optimal\dots

\paragraph{Notation.}

We have a graph $G = (V, E)$ and a length function $\ell : E \to \R_{\geq 0}$.
$d(s, t)$ denotes the length of the shortest path between $s$ and $t$,
($\ell(P) \defeq \sum_{e \in P}\ell(e)$). When $G$ is undirected denote the ball
of radius $r$ around $s \in V$ by
\[
  \ball{s, r} \defeq \Set{v \in V \mid d(s, v) \leq r}.
\]
When $G$ is directed, we have the in-ball and out-ball, defined respectively as
\begin{align*}
  \inball{s, r}  & \defeq \Set{v \in V \mid d(v, s) \leq r}  \\
  \outball{s, r} & \defeq \Set{v \in V \mid d(s, v) \leq r}.
\end{align*}

\section{Undirected LDDs}

\begin{definition}
  We have an undirected graph $G = (V, E)$ along with a nonnegative length
  function $\ell : E \to \R_{\geq 0}$. Let $\Delta \geq 0$ be a parameter (the
  ``diameter''). A \emph{low diameter decomposition} of $G$ is a partition of
  $V$ into sets $V_1, \dots, V_k$ so that
  \begin{itemize}
    \item For each $i$ and $s, t \in V_i$, $d(s, t) \leq \Delta$, i.e.\ the
          components have low diameter. Note the distance is with respect to the
          original graph; there is no guarantee on $d_{G[V_i]}(s, t)$.

    \item The probability that an edge is cut is small: $\Pr{\text{$u$ and
              $v$ are in different components}} \leq \apxO{\frac{\ell(u,v)}{\Delta}}$ for all edges $(u, v) \in E$.
  \end{itemize}
\end{definition}
Note that the partition must be randomized, since the second requirement uses
probability. It is far from obvious that such a decomposition even exists, let
alone is efficiently computable. In fact, there is an entire zoo of algorithms
solving this problem; we'll work through forming one below.

\subsection{Cutting Out Balls}

\def\Unif#1{\operatorname{Unif}\paren{#1}}
\def\Geom#1{\operatorname{Geom}\paren{#1}}
\def\Exp#1{\operatorname{Exp}\paren{#1}}

A simple type of algorithm is one that successively computes $V_i$'s, cutting
them out of the graph and continuing until the graph is empty. More
specifically, suppose we've determined $V_1$ through $V_{i-1}$. To form $V_i$,
we'll choose some vertex $s \in V \setminus (V_1 \cup \cdots \cup V_{i-1})$ and
some random radius $r$, and set $V_i = \ball{s, r}$. Here the ball is with
respect to distances in the remaining graph.

% badness-format skip
\begin{centeralgo}
  \textsc{GenericLDD}$(V, E, \ell, \Delta)$:\+
  \\  $i \gets 0$
  \\  while some vertex $s$ is not in $V_1 \cup \cdots \cup V_i$:\+
  \\    $i \gets i + 1$
  \\    sample $r \sim \,???$ \Comment{Some distribution...}
  \\    $V_i \gets \ball{s, r}$ \Comment{Dijkstra's in the remaining graph}\-
  \\  return $V_1, \dots, V_{i - 1}$
\end{centeralgo}
The above algorithm has the nice property of running in near-linear time, since
we never look at an edge more than twice (why?).

What distribution should we choose? We need every sampled radius to be at most
$\Delta/2$ with high probability, since then it's clear that every component has
diameter at most $\Delta$ (In fact, via this construction, the diameter is with
respect to $G[V_i]$ itself, i.e.\ it's \emph{strong}).

One natural option is to sample $r$ from $\Unif{0, \Delta/2}$, since that
guarantees that $r \leq \Delta/2$. What is the probability that some edge
$(u, v)$ is cut? Consider the first $V_i$ that contains any endpoint of the
edge; WLOG suppose $u \in V_i$. Let $s_i$ denote the center of the ball.
Conditioned on $u \in V_i$, $r \sim \Unif{d(s, u), \Delta}$. We know that
$d(s, v) \leq d(s, u) + \ell(u, v)$, but the proof starts to fall apart here.
What if $d(s, v) = d(s, u) + \ell(u, v) \geq \Delta/2$? The edge will always be
cut (assuming $u \in V_i$)!~\footnote{TODO: Show CKR-style scheme that fixes
this.}

Somehow, we'd want the distribution to play nice once we condition on
$u \in V_i$. An obvious candidate is a memoryless distribution, so let's try
sampling from $\Exp{\lambda}$ for some $\lambda$. Recall the pdf of this
distribution is $\lambda e^{-\lambda x}$, and so the probability that a sample
is below some threshold is
\begin{align*}
  \Pr_{X \sim \Exp{\lambda}}{X \leq t} = \int_0^t \lambda e^{-\lambda x}
                                       = \brack{-e^{-\lambda x}}_0^t
                                       = 1 - e^{-\lambda t}
\end{align*}
Now let's run through the argument with $(u, v)$ again. Conditioned on
$u \in V_i$, the chance that $v \notin V_i$ is at most the chance that
$r \leq d(s, u) + \ell(u, v)$. Since $r \geq d(s, u)$ and the distribution is
memoryless, this happens with probability
\[
  1 - e^{-\ell(u, v)\lambda} \leq \ell(u, v)\lambda,
\]
where the inequality is from a rearrangement of $1 - x \leq e^{-x}$. The
question now is how small can we make $\lambda$? We need $r \leq \Delta/2$ with
high probability, so setting $\lambda = \frac{2c\ln n}{\Delta}$ gives
\begin{align*}
  \Pr{r < \Delta/2} = 1 - e^{\frac{\Delta}{2}\frac{2c\ln n}{\Delta}}
                    = 1 - e^{-c \ln n}
                    = 1 - n^{-c}.
\end{align*}
We then have a cutting probability for $(u, v)$ of at most
\[
  \ell(u, v)\lambda = \bigO{\frac{\ell(u, v)\log n}{\Delta}}
                    = \apxO{\frac{\ell(u, v)}{\Delta}}.
\]

\begin{remark}
  Nothing here is special to graphs. The exact same algorithm works with any
  metric space (or even semi-metric space). The more general guarantee is
  \[
    \Pr{\text{$u$ and $v$ are in separate components}} \leq \bigO{\frac{d(u,
          v)\log n}{\Delta}}.
  \]
\end{remark}

\begin{remark}
  Consider the cycle graph with unit lengths and any $\Delta$. Any LDD of $C_n$
  must have components of length at least $\Delta$, so at least every $\Delta$th
  edge must be cut. Hence, if the probability that any edge is cut is at most
  $p$,
  \[
    n/\Delta \leq            \E{\text{\# of edges cut}}
             \leq            np
             \Longrightarrow p
             \geq            1/\Delta.
  \]
  This shows that we can't do better than $\ell(u, v)/\Delta$ in general. The
  $\log n$ factor is called the \emph{loss} of the LDD, and more complicated
  proofs show that it must be $\bigOmega{\log n}$.
\end{remark}

\section{Directed LDDs}

(This is based on my memory of~\cite{BNW25}.)

There is a directed analogue of LDDs, but it's less elegant. It helps to reframe
the goal of undirected LDDs as ``remove edges so that all connected components
have low diameter'' instead of ``partition the vertices into low diameter
pieces''. The directed analogue would then be ``remove edges so that all
\emph{strong} connected components have low diameter''.

From now on, for a set of edges $S$, let $G \setminus S$ denote the graph after
the removal of those edges.
\begin{definition}
  Let $G = (V, E)$ be a directed graph, $\ell : E \to \R_{\geq 0}$ a length
  function, and $\Delta$ a parameter. A \emph{directed low diameter
  decomposition} of $G$ is a (random) set $S \subseteq E$ such that
  \begin{itemize}
    \item For every SCC $C$ of $G \setminus S$ and $s, t \in C$,
          $d_G(s, t) \leq \Delta$.
    \item For every $(u, v) \in E$,
          $\Pr{(u, v) \in S} \leq \apxO{\frac{\ell(u, v)}{\Delta}}$.
  \end{itemize}
\end{definition}

The trickiness with algorithms here is that we lose a lot of the nice symmetry
we had before. In particular, the out-ball $\outball{s, r}$ for
$r \leq \Delta/2$ need not have diameter at most $\Delta$ (why?).

Let's recast the previous undirected LDD algorithm in the ``cutting edges''
lens. We computed $\ball{s, r}$, and then added the edges in $\cut{\ball{s, r}}$
to our cut set. After this, we immediately knew that no connected component of
$G$ contained vertices from both the ball and its complement. Something similar
occurs in the directed case. For any $U \subseteq V$, if the edges in
$\outcut{U}$ (or $\incut{U}$) are among our cut edges, then no SCC of the
remaining graph can contain vertices from both $U$ and $V \setminus U$. The
natural candidate for $U$ is a ball, but we lose the undirected symmetry that
guaranteed $u, v \in \ball{s, r} \Rightarrow d(u, v) \leq 2r$. To fix this,
we'll recursively decompose the ball as well.

% badness-format skip
\begin{centeralgo}
  \textsc{DirectedLDDAttempt}$(G = (V, E, \ell), \Delta)$:\+
  \\  if $d(s, t) \leq \Delta$ for every $s, t \in V$: \Comment{APSP}\+
  \\    return $\emptyset$\-
  \\  $i \gets 0$
  \\  while some vertex $s$ is not in $V_1 \cup \cdots \cup V_i$:\+
  \\    $i \gets i + 1$
  \\    sample $r \sim \,???$ \Comment{Some distribution...}
  \\    $V_i \gets \ball{s, r}$ \Comment{Dijkstra's in the remaining graph}
  \\     \Comment{After removing $S_0$, no $u \in V_i$ is strongly connected to a $v \notin V_i$}
  \\    $S_0 \gets S_0 \cup \outcut{V_i}$
  \\    $S_1 \gets S_1 \cup \textsc{DirectedLDDAttempt}(G[V_i], \Delta)$\-
  \\  return $S_0 \cup S_1$
\end{centeralgo}

Don't worry about optimizing runtime for the moment (running APSP makes this
$\bigOmega{mn}$ anyhow). If we pick the exponential distribution for $r$ like
before, it's not hard to see that the analysis essentially goes through the same
(exercise): each $(u, v) \in E$ is in $S_0$ with probability at most
$\apxO{\ell(u, v)/\Delta}$. By induction and a union bound, $(u, v) \in S$ with
probability $\apxO{q\ell(u, v)/\Delta}$, where $q$ is the number of calls that
$(u, v)$ is in. If each $V_i$ is at most a constant factor of $V$, then we have
$q = O(\log n) = \apxO{1}$.

Of course, we might get unlucky in our choice of $s$---if $s$ reaches more than
a constant factor of the graph, the recursion might get strung along for
$\poly(n)$ calls. Let's pick some $\beta < 1$ and say we want
$\sizeof{V_i} \leq \beta|V|$ for all $i$. If we can find an $s$ for which we're
confident that $\outball{s, r}$ is small, then we can make progress. We can pick
the distribution so that $r \leq \Delta/2$ with high probability, it suffices to
find an $s$ for which $\sizeof{\outball{s, \Delta/2}} \leq \beta \sizeof{V}$.
Failing this, we can find an $s$ for which the in-ball $\inball{s, \Delta/2}$
has small size--in this case we'd remove the in-cut rather than the out-cut.

For ease of notation, let $V_{<i} = \bigcup_{j<i}V_j$.
% badness-format skip
\begin{centeralgo}
  \textsc{DirectedLDD}$(G = (V, E, \ell), \Delta)$:\+
  \\  $i \gets 0, S_0 \gets \emptyset, S_1 \gets \emptyset$
  \\  repeat:\+
  \\    $i \gets i + 1$
  \\    sample $r \sim \Exp{\bigO{\log n / \Delta}}$
  \\    if there's some $s \notin V_{<i}$ with $\sizeof{ \outball{s, \Delta/2} } \leq
  \beta\sizeof{V}$:\+
  \\      $V_i \gets \outball{s, r} \setminus V_{< i}$
  \\      $S_0 \gets \outcut{V_i}$\-
  \\    else if there's some $s$ with $\sizeof{ \inball{s, \Delta/2} } \leq
  \beta\sizeof{V}$:\+
  \\      $V_i \gets \inball{s, r} \setminus V_{< i}$
  \\      $S_0 \gets \incut{V_i}$\-
  \\  else:\+
  \\    \Comment{Hope for the best\dots}
  \\    return $S_0 \cup S_1$\-
  \\  $S_1 \gets S_1 \cup \textsc{DirectedLDD}(G[V_i])$
\end{centeralgo}

It's clear that $\Pr{(u, v) \in S} = \bigO{\frac{\ell(u, v)\log^2 n}{\Delta}}$
with high probability, since $r \leq \Delta/2$ whp and so
$\sizeof{V_i} \leq \beta|V|$ for all $V_i$, meaning there are $\bigO{\log n}$
levels of recursion.

In an invocation, let $V_\text{rem}$ denote the vertices not in any $V_i$.
Taking stock, the algorithm partitions $V$ into $V_1, \dots, V_i, V_\text{rem}$;
these sets are guaranteed to be in distinct SCCs of $G \setminus S$ via $S_0$.
It recurses on $V_1$ through $V_i$; proving correctness comes down to proving
that $V_\text{rem}$ has (weak) diameter at most $\Delta$. Note that any
$v \in V_\text{rem}$ must have
$\sizeof{\outball{v, \Delta/2}}, \sizeof{\inball{v, \Delta/2}} \geq \beta\sizeof{V}$.
But if we choose any $\beta > 1/2$, for any $u, v \in V_\text{rem}$, the
out-ball of $u$ and the in-ball of $V$ must intersect! This implies that
$d(u, v) \leq \Delta/2 + \Delta/2 = \Delta$, so we indeed have our low-diameter
guarantee.

\paragraph{Running time.}

This algorithm takes $\apxO{mn}$ time. Each call can be implemented via $n$
calls to Dijkstra's algorithm; extracting balls once we know all distances is
trivial. One might wonder if we can achieve near-linear running time like we did
in undirected LDDs. The trickiness is in finding the ``in-heavy'' and
``out-heavy'' vertices which reach over $\beta|V|$ others--we turn to estimation
to efficiently detect these.

{\color{gray}
  \paragraph{Estimation.}

  (This is not interesting). The goal is to find which $v \in V$ have
  $\outball{v, \Delta/2} \geq \beta n$. If we can do this, we can reverse the
  graph to
  find the vertices with large in-balls. If $v$ reaches at least a $\beta$
  fraction of $V$ within $\Delta$, we'd expect it to reach some vertex in a
  random
  sample of $1/\beta$ vertices a constant fraction of the time. In particular,
  \begin{align*}
    \Pr{\text{$v$ reaches some vertex in $1/\beta$ ones}}
    =    1 - \Pr{\text{$v$ reaches none in the sample}}
    \geq 1 - (1 - \beta)^{1/\beta}
    \geq 1 - 1/e.
  \end{align*}
  Conversely, if $v$ reaches, say, at most a $2\beta/3$ fraction, running
  through
  the calculations gives a hitting probability of at most
  \[
    1 - (1 - 2\beta/3)^\beta \geq 1 - e^{-2/3}.
  \]
  Now we can simple run $O(\log n)$ experiments: in each one, we randomly
  sample
  $1/\beta$ vertices into $S$ and compute $d(\cdot, S)$ using $G^R$ to find
  $R_S = \Set{v \in V \mid d(v, S) \leq \Delta/2}$. Then, for each vertex, if
  it's
  in $R_S$ for at least a $p$ portion of the samples, we mark it as heavy.
  Here,
  $p$ is just some threshold between $1 - e^{-2/3}$ and $1 - e^{-1}$. Chernoff
  bounds imply that any vertex not marked as heavy reaches less than $\beta$
  of
  the graph, while vertices that are marked as heavy reach at least $2\beta/3$
  of
  the graph. Choosing $\beta = 4/5$ gives $2\beta/3 = 8/15 > 1/2$, so the
  pigeonhole argument for heavy vertices still applies.
}

\paragraph{Making it fast.}

We can figure out which vertices are (in/out)-heavy in $\bigO{m\log^2n}$ time
(via $\bigO{\log n}$ SSSP computations). This happens at the beginning of all
invocations of the recursive directed LDD algorithm. When cutting in or
out-balls $\inoutball{s, r}$, we run Dijkstra's algorithm only in the remaining
part of the graph. This ensures that we never repeatedly process edges, and the
resulting balls are still subsets of $\inoutball{s, \Delta/2}$ (with respect to
$G$'s distance) with high probability. All in all, at a recursive level we spend
$\bigO{m\log^2n}$ time over $\bigO{ \log n }$ levels, giving a runtime of
$\bigO{m\log^3 n}$.

Packaging this up as a lemma, we have the following.
\begin{lemma}
  Given $G = (V, E)$ and $\ell : E \to R_{\geq 0}$, we can compute a directed
  LDD $S$ of $G$ such that the probability of an edge $(u, v)$ being cut is
  $\bigO{ \frac{\ell(u, v)\log^2n}{\Delta} + \frac{1}{ n^c } }$. With high
  probability this takes $\bigO{m\log^3 n}$ time.
\end{lemma}
The $\bigO{\frac{1}{n^c}}$ is there to account for the failure probability of
the classification, in which case we can just return every edge.

\subsection{Later Improvements}

The above algorithm is roughly what the first near-linear weakly-polynomial SSSP
algorithm used~\cite{BNW25}. I do want to show a (from a certain perspective)
simplification of~\cite{Li25}, which, among other improvements, showed a nice
way to deal with heavy vertices without relying on pigeonhole arguments.

For some threshold $\beta < 1$ and $\Delta$, call a vertex $u$ \emph{out-heavy}
(resp. \emph{in-heavy}) if
$\sizeof{ \outball{u, \Delta/8} } \geq \beta \sizeof{V}$ (resp.
$\sizeof{\inball{u, \Delta/8}} \geq \beta \sizeof{V}$). Note that we have
tightened the threshold radius to $\Delta/8$ for reasons that will be apparent
soon. The analysis for light vertices survives.

Suppose there were no out-heavy vertices. Then we can repeatedly cut the
out-boundaries of out-balls and rest assured that subproblems go down by a
constant factor. Similarly, if there are no in-heavy vertices, then we can just
cut in-boundaries of in-balls. It turns out that avoiding mixed orientations of
the balls opens the door to better algorithms~\footnote{I don't fully see why
yet\dots} (see~\cite{Li25} for details; also see CKR's multiway cut rounding).
The challenge, then, is what to do in the case when there are both in and
out-heavy vertices. Let $S^-$ and $S^+$ denote the in and out heavy vertices
respectively.

\paragraph{Case 1.}

Suppose $d(S^-, S^+) > \Delta/2$. We \emph{uniformly} sample
$r \in [0, \Delta/2]$ and compute $U = B^+_r(S^-)$. Remove the edges in
$\outcut{U}$. A given edge $(u, v)$ is cut when
\[
  r \in [d(S^-, u), d(S^-, v)] \subseteq [d(S^-, u), d(S^-, u) + \ell(u, v)],
\]
which happens with probability at most $2\ell(u, v)/\Delta$. Note that
$U \supseteq S^-$ and $V \setminus U \supseteq S^+$, so no recursive call both
in and out heavy vertices. The recursive calls on are guaranteed to reduce the
problem by a constant factor\footnote{I'm being a little vague here, since this
is meant to be a general preprocessing step, but this is true for any algorithm
that cuts boundaries of balls. Furthermore, the constant factor reduction is
with respect to the size of the graph before this step, not the graph passed to
the calls.}, so this cut probability occurs at most $\bigO{\log n}$ times.

\paragraph{Case 2.}

Now suppose there is some $s^- \in S^-, s^+ \in S^+$ such that
$d(s^-, s^+) \leq \Delta/2$. Now sample $r$ uniformly in $[\Delta/8, \Delta/4]$
and compute $U^- = B^-_r(s^-)$ and $U^+ = B^+_r(s^+)$. Remove the edges in
$\incut{U^-} \cup \outcut{U^+}$. This splits $V$ into four regions: letting
$U = U^- \cup U^+$, we have
\begin{itemize}
  \item The outside $V \setminus U$
  \item The intersection $U^- \cap U_+$
  \item $U^+ \setminus U^-$
  \item $U^- \setminus U^+$
\end{itemize}
No SCC of $G \setminus (\incut{U^-} \cup \outcut{U^+})$ can contain vertices
from more than one of these regions (why?).
\begin{center}
  [Venn diagram with caption: No cycle can cross any of the boundaries.]
\end{center}

Exercise: every edge $(u, v)$ is cut with probability $\bigO{\ell(u, v)/\Delta}$
(the proof is very similar to case 1).

Moreover, every region except for $U^- \cap U^+$ is at most a constant factor of
the graph: each of them completely avoids either $U^-$ or $U^+$, which have
sizes of at least
$\sizeof{\inball{s^-, \Delta/8}}, \sizeof{\outball{s^+, \Delta/8}} \geq \beta\sizeof{V}$.
We can thus comfortably recurse on these components without blowing up depth.

Finally, there is no need to recurse on the intersection, since for any
$u, v \in U^- \cap U^+$,
\[
  d(u, v) \leq d(u, s^-) + d(s^-, s^+) + d(s^+, v)
          \leq \Delta/4 + \Delta/2 + \Delta/4
          =    \Delta.
\]

\begin{nudge}
  Think about why uniform distances worked here but not in our earlier LDD
  approaches.
\end{nudge}

\printbibliography
\end{document}
